\subsection{Finite-place cancellation in the relative determinant}
\label{geom:relative-finite-index}

Fix a word \(i\) at a translate \(u\), with the original arithmetic
record \(D,L\). Let \(M_{ui}\) be the base lattice, \(M^{\rm ram}_{ui}\)
its pullback to the descended vector space, \(N^0_{ui}\) the lattice
before the fractional-ideal shift, and \(N_{ui}\) the shifted lattice.
In the source these are the base \code{lat 0}, \code{MramL},
\code{NuLat}, and the descended \code{Nu}. The comparison
assumes the original commensurability, finite-support and local
length identifications, and the level-prime comparison
\code{adm26}. Write \(r_{ui}>0\) for the common rank of
\(N_{ui},N^0_{ui},M^{\rm ram}_{ui}\), and retain \(d^M_{ui}\) for the
rank of \(M_{ui}\). With the original full lattice and word rank
identifications, the two ranks have the relation
\(r_{ui}=d^M_{ui}(1-1/(T_{ui}-1))^f\), where \(T_{ui}>2\) is the
word level and \(f\) is the number of weight coordinates.
\src{src/Formalism/LocalRingCardinality/DataDIVDefs.lean}{12}{the pullback and shifted lattice}
\src{src/Formalism/LocalToGlobal/GlobalTheory/PlacesP.lean}{124}{the descended rank}
\src{src/Formalism/SharpRestriction/Divisibility/LocalQuadraticParts.lean}{621}{the level-prime comparison}.

For a finite place \(w\) of \(K_u\), let \(e_w\) denote its ramification
index and put \(\lambda_w=\log N\mathfrak p_w/[K_u:\mathbf Q]\).
Use \(g_{jw},F_{jw},E_{uijw}\) for the recorded gain, shift and
exponent correction at lift \(j\). In particular \(g_{jw}\) is zero
on the dropped primes. Set
\[
 s_{uiw}=\sum_j(g_{jw}+F_{jw}-E_{uijw}),\qquad
 \xi_{uiw}=e_ws_{uiw}-\lfloor e_ws_{uiw}\rfloor.
\]
Let \(\ell^{\rm ind}_{uiw}\) be the local length of
\(M^{\rm ram}_{ui}/N^0_{ui}\), divided by \(e_w\), and let
\(\tau_w\) be the original gain correction at \(w\).
The field-normalized slope formula gives the exact identity
\begin{align}
 \mu^N_{ui}-\mu^M_{ui}+E_{ui}-F_u^{\,0}
    &=G_u-\Diamond_{ui},\label{geom:finite-cancel}\\
 \Diamond_{ui}
   &=
   \frac{f\log(T_{ui}-1)}{T_{ui}-2}
     +2\sum_{w\in\mathrm{supp}}\lambda_w\xi_{uiw}
     +2\sum_{w\in\mathrm{supp}}\lambda_we_w
          \left(\frac{\ell^{\rm ind}_{uiw}}{r_{ui}}-\tau_w\right).
   \label{geom:finite-diamond}
\end{align}
Here \(E_{ui}=2\sum_{j,w}\lambda_we_w E_{uijw}\),
\(F_u^{\,0}=2\sum_{j,w}\lambda_we_w F_{jw}\), and
\(G_u=2\sum_w\lambda_we_w(\sum_jg_{jw}-\tau_w)\). All sums are over
the original finite place support and lift set.
\src{src/Formalism/LocalToGlobal/GlobalTheory/PlacesP.lean}{129}{the exact slope identity}
\src{src/Formalism/SharpRestriction/GlobalSchurBounds/TopDensityNGeSixteen.lean}{1489}{the factor of two and the three sums}.

Indeed the shift multiplies \(N^0_{ui}\) by
\(\prod_w\mathfrak p_w^{-\lfloor e_ws_{uiw}\rfloor}\), which adds
\(2\sum_w\lambda_w\lfloor e_ws_{uiw}\rfloor\) to the slope.
The finite index from \(N^0_{ui}\subseteq M^{\rm ram}_{ui}\)
subtracts \(2\sum_w\lambda_we_w\ell^{\rm ind}_{uiw}/r_{ui}\).
Both comparisons use the same descended vector space and its
determinant metric, so common different factors cancel in the
same-rank ratios. The comparison with \(M_{ui}\) uses the separate
level-prime slope correction in \eqref{geom:finite-diamond}. This
keeps the actual descended rank in the finite index denominator.
Substituting \(\lfloor e_ws_{uiw}\rfloor=e_ws_{uiw}-\xi_{uiw}\)
cancels \(E_{ui}\) and \(F_u^{\,0}\) with the indicated signs.
\src{src/Formalism/SharpRestriction/AdmissibleParameterJoins/SumLogPrimeFactorsLe.lean}{358}{the fractional-ideal norm}
\src{src/Formalism/SizeWindowBrackets/ChainSlopeWindows/ExistsNsmulMem.lean}{238}{the same-rank index ratio}.

\paragraph{The sign of the remaining local index.}
Use the original \code{LocLe} comparison off the word-level
prime and the integral values of \(\tau_w\ge0\) that enter its
construction. In the DVR \(\mathcal O_{K_u,\mathfrak p_w}\), these
give
\((N^0_{ui})_w\subseteq\mathfrak p_w^{e_w\tau_w}
(M^{\rm ram}_{ui})_w\). The localized pullback lattice is free of
rank \(r_{ui}\). The quotient by \((N^0_{ui})_w\) therefore maps onto
a direct sum of \(r_{ui}\) copies of
\(\mathcal O_{K_u,\mathfrak p_w}/\mathfrak p_w^{e_w\tau_w}\),
of total length \(r_{ui}e_w\tau_w\). Dividing by \(r_{ui}e_w\)
gives \(\ell^{\rm ind}_{uiw}/r_{ui}\ge\tau_w\).
At the word level \(\tau_w=0\) and the length is nonnegative.
Writing \(c_{T_{ui}}=f\log(T_{ui}-1)/(T_{ui}-2)\), we obtain
\(\Diamond_{ui}\ge c_{T_{ui}}\ge0\) for the original \(T_{ui}>2\).
\src{src/Formalism/LocalRingCardinality/DataDIVDefs.lean}{32}{local inclusion}
\src{src/Formalism/LocalToGlobal/GlobalTheory/PlacesP.lean}{297}{the inputs for the original place comparison}
The original proof gives
\src{src/Formalism/LocalToGlobal/GlobalTheory/PlacesP.lean}{521}{integrality of the correction}
and uses it in \code{PlacesP.lean}, lines 729--736.

With the same \(P_{ui}=\mu^N_{ui}-\mu^M_{ui}+E_{ui}
-F_u^{\,0}-\mathcal M_u\) and \(\ell_{ui}\) as above, the lower
quotient slope bound now gives
\[
 G_u-\mathcal M_u=P_{ui}+\Diamond_{ui},\qquad
 P_{ui}+c_{T_{ui}}\le G_u-\mathcal M_u\le P_{ui}+\ell_{ui}.
\]
The term \(\Diamond_{ui}\) enters the signed discrepancy with a
positive sign. For any one chosen word \(i\) at \(u\) one has
\(D_u=(\mathcal A_u-G_u)+P_{ui}+\Diamond_{ui}\).
This permits a test with one actual word per translate and does not
require a separate maximum-slope variation bound. It still requires
an upper estimate for the actual relative determinant term.
\src{src/Formalism/SizeWindowBrackets/QuotientSlopeBoundsWith.lean}{19}{the word loss and lower quotient slope}.

\paragraph{Integral denominator contributions.}
Rounding occurs only at places with \(e_ws_{uiw}\notin\mathbf Z\).
Consider a place in the support above a rational prime
\(p>X''\), with \(p\) different from the word level, the two generator
primes and the discriminant prime. Assume \(p\) is dropped and is
outside the generator types. The original definitions then give
\(g_{jw}=F_{jw}=0\) and \(E_{uijw}=v_p(D_{\rm Gam})\). The last
quantity can be nonzero, but \(e_ws_{uiw}=-e_wm v_p(D_{\rm Gam})\)
is an integer and contributes no rounding. If \(\tau_w=0\), the
original local equality at a dropped, non-level prime also gives
zero index. This removes these denominator places from the defect
even when they remain in the support of the individual exponents.
\src{src/Formalism/LocalRingCardinality/DataForced.lean}{126}{the prime branches for the shift}
\src{src/Formalism/LocalRingCardinality/DataErr.lean}{159}{the denominator valuation in the exponent}
\src{src/Formalism/LocalRingCardinality/DataDIVDefs.lean}{48}{local equality at a dropped prime}.

The Lean module \code{RelativeFinitePart.lean} checks
\eqref{geom:finite-cancel}, the support bound for the rounding loss,
and the vanishing index at a dropped non-level prime with zero
correction. It derives the nonnegative local excess from the original
local containment, the recorded local length, and the two correction
conditions above, using the existing DVR index theorem. Its resulting
signed slope bound retains \(c_{T_{ui}}\).
The DVR inequality used here is proved in
\code{DvIdxCalcMain.lean}, line 95.
